\cotefile{37}
\scanpage{52}
\noindent SGA 5 \(\overline{\underline{\mathrm{VII}}}\)\quad th de Lefschetz faible sans supposer \(X\) ni \(Y\) lisses (seulement \(X-Y\) lisse \ldots)

\noindent SGA 5 \(\overline{\underline{\unsure{\mathrm{V}|\mathrm{VI}}}}\)\quad \(\mathbb{Q}_{\ell}\)-faisceau : \unsure{signaler|rajouter} définition raisonnable en termes de \uline{champ} \uline{associé}

\noindent SGA 5. \tikz[baseline=(cancel.base)]{
\node[inner sep=1pt] (cancel) {\(\deleted{\overline{\underline{\unsure{\mathrm{IV}|\mathrm{V}|\mathrm{VI}}}}}\)};
\node[anchor=north,inner sep=1pt] (replacement) at ([yshift=-5pt]cancel.south) {\inserted{\(\overline{\underline{\mathrm{VI}}}\)}};
\draw[thin] ([xshift=-3pt,yshift=1pt]replacement.north west) .. controls +(-2pt,-16pt) and +(-10pt,-10pt) .. ([yshift=-4pt]replacement.south) .. controls +(16pt,-3pt) and +(2pt,-12pt) .. ([xshift=5pt,yshift=4pt]cancel.south east);
}\quad Reprendre en langage de \(\mathbb{Z}_{\ell}\)-faisceaux ou\textsuperscript{\inserted{\unsure{*|+}}} \(\mathbb{Q}_{\ell}\)-faisceaux les notions essentielles dont il a été question dans SGA 4 (, Ajouter nouveau N\textsuperscript{\uline{o}}).

\noindent SGA 5 \(\overline{\underline{\mathrm{IV}}}\)\quad Essayer de faire dans le cas relativement lisse ttes les constructions essentielles

\figure{\begin{tikzpicture}[x=\linewidth,y=1cm]
\draw[thin] (0,0) .. controls (.23,.12) and (.49,.08) .. (.62,.06)
 .. controls (.77,.05) and (.90,-.03) .. (.96,.02);
\end{tikzpicture}}

\noindent \(\overline{\underline{\mathrm{VI}}}\)\quad \unsure{Systèmes|Système} projectifs \inserted{\tikz[baseline=(adiques.base)]{
\node[inner sep=0pt] (adiques) {adiques};
\draw[thin] ([xshift=-.35em,yshift=-1.3ex]adiques.south west) .. controls +(90:1.2ex) and +(-.8em,-.25ex) .. ([yshift=-.2ex]adiques.south) .. controls +(.8em,.1ex) and +(0,-.8ex) .. ([xshift=.6em,yshift=.5ex]adiques.east);
}} dans les catégories \unsure{dérivées} \ldots\ (ou dans thèse)
% The following are faint isolated marks, not assigned a verbal reading.
\figure{\begin{tikzpicture}[x=\linewidth,y=1cm]
\path[use as bounding box] (0,-.24) rectangle (1,.20);
\node[inner sep=0pt] at (.35,.08) {\unsure{\tikz[x=1cm,y=1cm]{\draw[gray,thin] (0,0) -- (.08,.04) -- (.14,-.01) .. controls (.17,.06) and (.17,.14) .. (.24,.20);}}};
\node[inner sep=0pt] at (.45,-.08) {\unsure{\tikz[x=1cm,y=1cm]{\draw[gray,thin] (0,0) -- (.05,.08) -- (.10,.02) -- (.12,.17);\draw[gray,thin] (.10,.02) -- (.18,-.03);\draw[gray,thin] (.08,-.10) -- (.14,-.07);}}};
\end{tikzpicture}}
% Two oblique dots at the right, below the parenthetical note; their function is uncertain.
\margin{\unsure{\tikz[baseline=-.3ex]{\fill (0,0) circle[radius=.35pt];\fill (.07,.10) circle[radius=.35pt];}}}

\noindent SGA 5 \(\mathrm{V}\)\quad dans introduction du \unsure{livre|Hom} ou de l'expo\unsure{sé} signaler l'ennui des coeff. trop particuliers \((\mathbb{Z}_{\ell},\,\mathbb{Z}/\ell^{\nu}\mathbb{Z}\ -\,-)\)

\figure{\begin{tikzpicture}[x=\linewidth,y=1cm]
\draw[thin] (0,.03) .. controls (.12,.11) and (.19,.04) .. (.29,-.02)
 .. controls (.39,-.08) and (.45,.02) .. (.57,.01)
 .. controls (.74,.02) and (.85,-.11) .. (.91,-.06);
\end{tikzpicture}}

\noindent \(\mathrm{I}\)\quad \S\ 5\qquad Donner, si \(i\unsure{{:}|{<}|{\in}}\unsure{\eta|y}\longrightarrow X\) inclusion du p\textsuperscript{t} générique, et \(F_{\eta}\) Module loc. \unsure{libre} sur \(\eta\)
% The recurring eta/y ambiguity is marked at its first occurrence above.
\[
 D\bigl(i_{*}(F_{\eta})\bigr)=i_{*}\bigl(D(F_{\eta})\bigr)
\]
% archive sheet number: 51
\clearpage
